{\color{ink}\bfseries 例 2-10 \(\cdot \) 第1问}\par\vspace{1bp}只看循环控制，\(n=\allowbreak{}0\)时unsigned的n\(-\)1按模\(2^{32}\)变为FFFFFFFFH，\(i\le \)该上界恒真，i自增也会绕回；若i、n均为int，初次检查\(0\le -1\)为假，直接返回1，保留unsigned时可将条件改为\(i<\allowbreak{}n\)。原程序第31轮的\(power*=\allowbreak{}2\)触发有符号溢出，C11/\allowbreak{}C17无法保证实际表现为死循环，也不规定f1(31)的返回值。仅在额外约定32位补码运算按模\(2^{32}\)绕回时，f1(31)的sum才变成FFFFFFFFH并解释为\(-\)1；\(f1(30)=\allowbreak{}2^{31}-1\)，无溢出且精确返回的最大n为30。
